\documentclass[11pt]{article}
\usepackage[margin=1in]{geometry}
\usepackage[T1]{fontenc}
\usepackage{lmodern,microtype,amsmath,amssymb,amsthm,mathtools,booktabs}
\usepackage[colorlinks=true,allcolors=blue]{hyperref}
\usepackage{enumitem}
\setlist{nosep}
\newtheorem{theorem}{Theorem}
\newtheorem{lemma}[theorem]{Lemma}
\newtheorem{corollary}[theorem]{Corollary}
\newtheorem{remark}[theorem]{Remark}
\newcommand{\E}{\mathbb E}
\newcommand{\Prob}{\mathbb P}
\newcommand{\R}{\mathbb R}
\newcommand{\1}{\mathbf 1}
\newcommand{\pos}[1]{\left(#1\right)_+}
\newcommand{\TV}{\operatorname{TV}}
\newcommand{\KL}{\operatorname{KL}}
\newcommand{\Reg}{\operatorname{Reg}}
\newcommand{\iid}{\mathrm{iid}}
\newcommand{\prodclass}{\mathrm{prod}}
\newcommand{\calV}{\mathcal V}
\title{Separating iid and time-varying replication\\in online learning}
\author{Jitendra Prajapati}
\date{26 October 2026}
\begin{document}
\maketitle
\begin{abstract}
We compare two requirements for a full-information experts learner: its
entire action path must agree on independent input streams drawn either
from every iid law, or from every time-varying product law, when the runs
share their complete internal randomness. In both cases regret is
worst-case expected regret on every deterministic input, against one
fixed expert. For $n=2^m$ experts and $T=mL$ rounds, we give a finite
product-replication lower bound of order
$\min\{T,\sqrt T\,m^{3/2}/\rho\}$ at fixed replication error $\rho<1$.
A self-contained iid-replicable construction has regret
$O(\rho^{-2}\sqrt T\log(2n)[1+\log\log(2T)]^2)$.
At $T=m^3,n=2^m$, product-law replication forces linear regret while
iid replication permits sublinear regret. At $T=m2^m,n=2^m$, the resulting
multiplicative gap diverges even though $n<T$.
The lower bound allows arbitrary shared randomness, hidden state and
rare exceptional branches. These are growing-expert results; they do
not establish an impossibility of no-regret learning for fixed $n$,
an optimal general regret rate, or a practical failure of trained models.
\end{abstract}

\section{Model, statements and scope}

There are $n$ experts and $T$ rounds. Before seeing
$c_t\in[0,1]^n$, a causal learner chooses an index
$I_t=I_t(c_{<t},R)$, where its full random tape $R$ is independent of
the inputs. For a deterministic input $c=(c_1,\ldots,c_T)$ define
\[
 \Reg_T(c,R)=\sum_{t=1}^T c_t(I_t)-\min_{i\in[n]}\sum_{t=1}^T c_t(i),
 \qquad
 \mathcal R_T=\sup_{c\in([0,1]^n)^T}\E_R\Reg_T(c,R).
\]
The expectation always includes the \emph{entire original tape};
we impose no regret condition on individual seeds.

For a fixed law $\mathcal D_1\otimes\cdots\otimes\mathcal D_T$ on the
cost vectors, draw $c,c'$ independently and use the same $R$.
The learner is \emph{product $\rho$-replicable} if
\begin{equation}\label{eq:rep}
 \Prob\bigl[(I_t(c_{<t},R))_{t=1}^T
          \ne(I_t(c'_{<t},R))_{t=1}^T\bigr]\le\rho
\end{equation}
for every such fixed product law. It is \emph{iid $\rho$-replicable}
if the condition is required only when all $\mathcal D_t$ are the
same unknown distribution. Coordinates within each cost vector may be
arbitrarily dependent. The distributions are fixed independently of
the realized inputs and tape; no adaptive-input assertion is made.

These are the two replication notions considered by
Ahmadi, Bhandari and Blum~\cite{ABB}.
Let $\calV_{\iid}(T,n,\rho)$ and $\calV_{\prodclass}(T,n,\rho)$ be the
infima of $\mathcal R_T$ over the respective classes of learners.
Thus the two values have the \emph{same} regret criterion;
only the distribution class in~\eqref{eq:rep} changes.

\begin{theorem}[Finite product-replication lower bound]\label{thm:lower}
Let $m,L$ be positive integers, $T=mL$, $n=2^m$, and $0<\rho<1$.
Every product $\rho$-replicable learner satisfies
\begin{equation}\label{eq:lower}
 \mathcal R_T\ge
 \min\left\{\frac{(1-\rho)T}{32},
 \frac{\sqrt3(1-\rho)^3\sqrt T\,m^{3/2}}{768\rho}\right\}.
\end{equation}
The bound permits arbitrary correlations throughout the shared tape,
arbitrary internal state, and a global rare branch.
\end{theorem}

\begin{theorem}[An iid-replicable upper bound]\label{thm:upper}
For integers $n,T\ge2$ and $0<\rho<1$, put
\[
 J=\left\lceil\log_2\max\{1,\log T\}\right\rceil,\qquad B=J+2,
\]
\[
 U=2\sqrt T+\frac{16B^2\sqrt T\log n}{\rho}
          +\frac{\rho\sqrt T}{8},\qquad
 A=U+3\sqrt{T\log(2n)},\qquad \kappa=\frac{8A}{\rho}.
\]
There is a causal iid $\rho$-replicable learner with
\begin{equation}\label{eq:upper}
 \mathcal R_T\le
 \min\left\{T,\ \kappa+1+\sqrt{\frac{T\log n}{2}}\right\}
 =O\!\left(\min\left\{T,
 \frac{\sqrt T\log(2n)[1+\log\log(2T)]^2}{\rho^2}\right\}\right).
\end{equation}
It uses $O(nT)$ ideal real-arithmetic operations and
$O(n+\log\log T)$ working storage, with exponential sampling and
elementary real functions as primitives. This is not a random-bit
or floating-point complexity claim.
\end{theorem}

\begin{corollary}[Two separation regimes]\label{cor:separation}
Fix $0<\rho<1$.
\begin{enumerate}
\item At $T=m^3,n=2^m$, as $m\to\infty$,
\[
 \calV_{\prodclass}(m^3,2^m,\rho)=\Theta_\rho(m^3),\qquad
 \calV_{\iid}(m^3,2^m,\rho)
 =O_\rho\bigl(m^{5/2}[1+\log\log m]^2\bigr)=o(m^3).
\]
\item At $T=m2^m,n=2^m$, for integers $m\ge2$,
\[
 \calV_{\prodclass}(T,n,\rho)=\Omega_\rho(\sqrt T\,m^{3/2}),\qquad
 \calV_{\iid}(T,n,\rho)
 =O_\rho\bigl(\sqrt T\,m[1+\log m]^2\bigr).
\]
Consequently $\calV_{\prodclass}/\calV_{\iid}$ tends to infinity,
at least as a constant times $\sqrt m/[1+\log m]^2$, although $n<T$.
\end{enumerate}
\end{corollary}

\begin{proof}
Substitute $L=m^2$ and $L=2^m$, respectively, in~\eqref{eq:lower}
and apply~\eqref{eq:upper}. In the first regime the two terms in
the lower bound are constant multiples of $T$, and the trivial regret
upper bound is $T$. In the second regime the square-root term is
eventually the smaller term and $\log\log T=\Theta(\log m)$.
The denominator in the ratio is positive: for any learner, choose a
least-probable first-round expert to have cost zero, give all other
experts cost one on that round, and give all later costs zero.
The resulting expected regret is at least $1-1/n>0$.
Both displayed comparison rates in the second regime are sublinear;
that lower bound does not establish an impossibility of no-regret
learning there.
\end{proof}

\paragraph{Relation to existing results.}
The main point is the stronger product-replication lower bound and
its qualitative consequence. The upper construction supplies a
self-contained sufficient iid comparison; it is not claimed to improve
the iid rate stated in~\cite{ABB}. The ingredients are classical Gibbs
stability, correlated sampling, martingale maximal bounds and a
regret-triggered fallback. Hopkins et al.~\cite{HIKLY} already study
adaptive sequential composition, including a sample cutoff followed by
a nonreplicable fallback. Here the different constraint is that actions
must be committed at prescribed rounds of a fixed input stream and the
entire clock-time action path must agree. Realizable PAC
learning~\cite{LMPS} instead concerns a final classifier.
The comparison is not a claim that these literatures are unrelated,
nor a claim of first discovery. Section~\ref{sec:discussion} specifies
the remaining boundaries.

\section{A causal phase-cube lower bound}

\subsection{One legal product family and one fixed comparator}

Index experts by $b\in\{-1,+1\}^m$. Partition the horizon into $m$
public phases of length $L$. Fix
$\theta=(\theta_1,\ldots,\theta_m)\in[-\delta,\delta]^m$,
where $0<\delta\le1/2$. Let the signs $X_{j,s}$ be independent,
with mean $\theta_j$, and give expert $b$ the cost
\begin{equation}\label{eq:cube}
 c_{j,s}(b)=\frac{1-b_jX_{j,s}}2.
\end{equation}
For each fixed $\theta$, this is a valid temporal product law.
Let $A_{j,s}\in\{-1,+1\}$ be the $j$th coordinate of the selected
expert, and write $S_j=\sum_{s=1}^L X_{j,s}$.

The realized best \emph{single} cube expert chooses
$b_j=\operatorname{sign}(S_j)$ simultaneously for all $j$.
Thus its expected-regret comparison is exactly
\begin{equation}\label{eq:regexact}
 \E\Reg_T=\frac12\sum_{j=1}^m
   \left(\E|S_j|-\sum_{s=1}^L\E[A_{j,s}X_{j,s}]\right).
\end{equation}
Causality gives $\E[A_{j,s}X_{j,s}]=\theta_j\E A_{j,s}$,
and $\E|S_j|\ge L|\theta_j|$. Introduce independent shared
uniform indices $K_j\in\{1,\ldots,L\}$, independent of everything
else. They are analysis-only coins, not inputs to the learner.
Equation~\eqref{eq:regexact} implies
\begin{equation}\label{eq:meanoracle}
 \E\Reg_T\ge
 L\sum_{j=1}^m|\theta_j|
       \Prob[A_{j,K_j}\ne\operatorname{sign}(\theta_j)].
\end{equation}
The right side is nonnegative after averaging over the full tape
and input law; we have not discarded a negative-regret seed.

For analysis, draw the coordinates $\theta_j$ independently and
uniformly from $[-\delta,\delta]$. We average guarantees valid at
\emph{each fixed} $\theta$. The resulting mixture is not itself
asserted to be a temporal product law. Include the $K_j$ in the
augmented shared tape $\bar R$ with its original distribution.

Define the phase probe $B_j=(A_{j,K_j}+1)/2$. Each probe is determined
before any future phase arrives. Equal entire action paths imply
equal probe paths. Before phase $j$, let $H$ be the complete earlier
raw sign blocks. For fixed $H=h,\bar R=r$, define
\[
 q_{h,r}(\vartheta)
 =\Prob[B_j=1\mid H=h,\bar R=r,\theta_j=\vartheta].
\]
The new coordinate $\theta_j$ is independent of the past and tape.
All internal state determined by $h,r$ is permitted.

Let $I_-=[-\delta,-\delta/2]$ and $I_+=[\delta/2,\delta]$, and put
\[
 \Delta_{h,r}=\E_{\vartheta\sim\mathrm{Unif}(I_+)}q_{h,r}(\vartheta)
             -\E_{\vartheta\sim\mathrm{Unif}(I_-)}q_{h,r}(\vartheta),
 \qquad v_j=\E_{H,\bar R}\Delta_{H,\bar R}.
\]
These are unconditional past-history averages, not averages
conditioned on earlier agreement. If
\[
 e_j=\E\bigl[|\theta_j|
    \1\{A_{j,K_j}\ne\operatorname{sign}(\theta_j)\}\bigr],
\]
then~\eqref{eq:meanoracle} gives
$L\sum_j e_j\le\mathcal R_T$. Each outer bin has prior probability $1/4$
and $|\theta_j|\ge\delta/2$ there. Hence
$e_j\ge(\delta/8)(1-v_j)$ and
\begin{equation}\label{eq:response}
 \frac1m\sum_{j=1}^m v_j\ge1-\frac{8\mathcal R_T}{\delta T}.
\end{equation}

\subsection{A one-phase transition inequality}

\begin{lemma}\label{lem:transition}
For every fixed past $h$ and augmented tape $r$,
\begin{equation}\label{eq:coarea}
 \E_{\vartheta\sim\mathrm{Unif}[-\delta,\delta]}
  q_{h,r}(\vartheta)(1-q_{h,r}(\vartheta))
 \ge\frac{\sqrt3}{12\delta\sqrt L}\,\Delta_{h,r}^2.
\end{equation}
\end{lemma}

\begin{proof}
Suppress $h,r$. The probe is a deterministic Boolean function of
the $L$ fresh signs when the past and full tape are fixed.
Its expectation $q(\vartheta)$ is a polynomial. The score identity
and Cauchy--Schwarz give
\[
 |q'(\vartheta)|
 \le\sqrt{\frac{Lq(\vartheta)(1-q(\vartheta))}{1-\vartheta^2}}
 \le\sqrt{L/3}=:M_0 ,
\]
since $|\vartheta|\le1/2$. This allows all $L$ signs, even if
the causal probe only uses a prefix.

Set $\Psi(u)=u^2/2-u^3/3$. For $a\in I_-$ and $b\in I_+$,
\[
 \begin{split}
 M_0\int_{-\delta}^{\delta}q(1-q)\,d\vartheta
 &\ge\int_a^b q(1-q)|q'|\,d\vartheta\\
 &\ge|\Psi(q(b))-\Psi(q(a))|
 \ge\frac{|q(b)-q(a)|^2}{6}.
 \end{split}
\]
The last inequality follows because, on any subinterval of $[0,1]$
of length $z$, the integral of $u(1-u)$ is at least
$z^2/2-z^3/3\ge z^2/6$. No monotonicity of $q$ is assumed.
Average over independent uniform $a,b$ in the two outer bins,
apply Jensen, and divide by $2\delta M_0$.
\end{proof}

\subsection{First mismatches with the entire original tape}

Let $D_{j-1}$ be the event that the previous probe paths disagree
between the paired runs, and let $d_{j-1}=\Prob(D_{j-1})$ under
the parameter prior, paired samples conditional on the parameter
and the full shared $\bar R$. We have $d_{j-1}\le\rho$, because
each fixed-parameter whole-probe failure is bounded by~\eqref{eq:rep}.

Conditional on earlier histories $H=h,H'=h'$ and tape $\bar R=r$,
the current raw sign blocks are independent given the fresh
$\theta_j$. Their probe probabilities are
$q=q_{h,r}(\theta_j)$ and $q'=q_{h',r}(\theta_j)$, so
\[
 \Prob[B_j\ne B'_j\mid h,h',r,\theta_j]
  =q+q'-2qq'\ge q(1-q).
\]
The event $D_{j-1}$ is determined by the earlier histories and tape,
and is independent of the new parameter and current blocks.
With $a_0=\sqrt3/12$, Lemma~\ref{lem:transition} therefore gives
\[
 d_j-d_{j-1}\ge
 \frac{a_0}{\delta\sqrt L}
 \E\left[\1_{D_{j-1}^c}\Delta_{H,\bar R}^2\right].
\]
The earlier histories may be correlated after their old parameters
are integrated out. No independence of $H$ and $H'$ is used.

Since $\Delta\le1$, deleting $D_{j-1}$ from its unconditional
expectation costs at most $d_{j-1}$. Cauchy--Schwarz implies
\[
 \E[\1_{D_{j-1}^c}\Delta^2]
 \ge\frac{\pos{v_j-d_{j-1}}^2}{1-d_{j-1}}
 \ge\pos{v_j-\rho}^2.
\]
Sum the disjoint first-mismatch increments, then use convexity and
\eqref{eq:response}, to obtain
\begin{equation}\label{eq:master}
 \rho\ge\frac{a_0m}{\delta\sqrt L}
       \pos{1-\rho-\frac{8\mathcal R_T}{\delta T}}^2.
\end{equation}
This is not a union-bound lower estimate on repeated mismatches,
nor an assumption of error $\rho/m$ in each phase. Accuracy was
not conditioned on previous agreement.

\begin{proof}[Proof of Theorem~\ref{thm:lower}]
Choose
\[
 \delta=\min\left\{\frac12,\,
      \frac{a_0m(1-\rho)^2}{4\rho\sqrt L}\right\}.
\]
Then $\sqrt{\rho\delta\sqrt L/(a_0m)}\le(1-\rho)/2$.
Rearranging~\eqref{eq:master}, including the case where its positive
part is zero, gives $\mathcal R_T\ge\delta T(1-\rho)/16$.
Substitution yields~\eqref{eq:lower}. Every parameter vector in this
argument defines a legal fixed product law; averaging over the prior
does not change the learner's required input model.
\end{proof}

\paragraph{Embedding.}
For more than $2^m$ experts, duplicate cube experts and map the chosen
index to its cube vector. Equal index paths still imply equal probe
paths. For $1\le m\le T$ and a horizon not divisible by $m$, use
$m\lfloor T/m\rfloor$ active rounds and append zero costs.
The lower formula then uses the active horizon. Neither embedding
adds restrictions to the original tape.

\section{A self-contained iid upper bound}

\subsection{Gibbs stability and a shared sampler}

For a cumulative cost vector $C$, write
$p_i(C)=e^{-\eta C_i}/Z(C)$.

\begin{lemma}\label{lem:gibbs}
Let $\eta>0$, let $P$ be a positive integer, and let $C,C'$ be
independent same-law sums of $P$ independent random vectors in $[0,1]^n$, not necessarily identically distributed.
Then
\[
 \E\KL(p(C)\Vert p(C'))\le\frac{\eta^2P}{2},
 \qquad
 \E\TV(p(C),p(C'))\le\frac{\eta\sqrt P}{2}.
\]
Using the same independent unit exponentials $E_i$ to sample
$\arg\min_i E_i/p_i(C)$ and $\arg\min_i E_i/p_i(C')$
gives mismatch probability at most $\eta\sqrt P$.
For iid summands of mean $\mu$,
\begin{equation}\label{eq:gibbsrisk}
 \E\langle p(C),\mu\rangle-\min_i\mu_i
 \le\frac{\log n}{\eta P}+\frac{\eta}{2}.
\end{equation}
\end{lemma}

\begin{proof}
Changing one summand changes each coordinate of $C$ by at most one.
Interpolate along $C(u)=C+ud$ for $0\le u\le1$ and
$d\in[-1,1]^n$, writing $p_i(u)=p_i(C(u))$. Differentiation gives
$p_i'(u)=-\eta p_i(u)(d_i-\E_{p(u)}d)$.
A random variable in an interval of length two has mean absolute
deviation at most one, so the resulting total variation is at most
$\eta/2$.

Replace summand $X_s$ by an independent copy, writing $C^{(s)}$.
The new sum has the original distribution and is independent of $X_s$.
Since coordinates of $X_s$ are in $[0,1]$,
\[
 \E\langle p(C),\E X_s-X_s\rangle
 =\E\langle p(C^{(s)})-p(C),X_s\rangle
 \le\eta/2 .
\]
The exact identity
\[
 \KL(p(C)\Vert p(C'))=
 \eta\langle p(C),C'-C\rangle+\log Z(C')-\log Z(C)
\]
has cancelling expected log-partition terms. Summing the preceding
inequality proves the KL bound. Pinsker and Jensen prove the
total-variation bound.

For deterministic positive distributions $p,q$, let $v=\TV(p,q)$.
The probability that $i$ wins both shared exponential races is
\[
 \left[\sum_j\max(p_j/p_i,q_j/q_i)\right]^{-1}
 \ge\frac{\min(p_i,q_i)}{1+v}.
\]
Summing proves mismatch at most $2v/(1+v)\le2v$.

Finally Gibbs' variational inequality gives
$\langle p(C),C\rangle\le C_i+\log n/\eta$ for every $i$.
Combine this with the replace-one estimate, sum over $s$, and compare
with an index minimizing $\mu_i$ to obtain~\eqref{eq:gibbsrisk}.
\end{proof}

\subsection{A pure learner on a deterministic calendar}

Let $J,B$ be as in Theorem~\ref{thm:upper}. Set
\[
 P_j=\left\lceil T^{\,1-2^{-(j+1)}}\right\rceil
       \quad(0\le j\le J),
\]
delete repeated values, and append $T$ if it is not already present.
Play expert 1 through $P_0$. After each retained prefix $P<T$,
use a fresh shared exponential race to draw from the Gibbs distribution
of the entire prefix with
\[
 \eta_P=\frac{\rho}{8B\sqrt P}.
\]
Keep that expert for the next calendar block. There are at most $B$
races. All race arrays belong to the original tape and are independent
of inputs.

Every subsequent block length $\ell$ satisfies
\begin{equation}\label{eq:calendar}
 \ell/\sqrt P\le2\sqrt T,\qquad P_0\le2\sqrt T.
\end{equation}
Indeed, successive unrounded prefixes satisfy
$x_{j+1}=\sqrt{Tx_j}$, so a nonfinal successor obeys
$P_{\rm next}\le\sqrt{TP}+1$. If repetitions were deleted, use the
last index having the repeated value. For the appended final block,
$2^{J+1}\ge2\log T$ gives $P_J\ge Te^{-1/2}$, and
$T/\sqrt{P_J}\le e^{1/4}\sqrt T<2\sqrt T$.

Lemma~\ref{lem:gibbs} and a union bound over this deterministic
calendar give pure entire-path mismatch probability at most $\rho/8$.
Under iid costs of mean $\mu$, with $\mu_*=\min_i\mu_i$, the
nonnegative random pseudo-regret
$W=\sum_t(\mu_{I_t}-\mu_*)$ obeys, by~\eqref{eq:gibbsrisk}
and~\eqref{eq:calendar},
\begin{equation}\label{eq:pseudoreg}
 \E W\le
 2\sqrt T+\frac{16B^2\sqrt T\log n}{\rho}
          +\frac{\rho\sqrt T}{8}=U .
\end{equation}
The race winner is independent of all subsequent block costs.
Using it throughout the block does not change this expectation.

\subsection{Controlling the maximum prefix regret}

For the pure learner define $R_0=0$ and
\[
 R_t=\sum_{s\le t}c_s(I_s)-\min_i\sum_{s\le t}c_s(i).
\]
Let
\[
 M_t=\sum_{s\le t}(c_s(I_s)-\mu_{I_s}),\qquad
 N_{i,t}=\sum_{s\le t}(c_s(i)-\mu_i).
\]
In the filtration containing the entire independent tape and observed
costs, these are martingales with increments of conditional mean zero
and range length at most one. Hoeffding's exponential supermartingale
and its maximal inequality yield
\[
 \Prob\!\left[\max_{t\le T}M_t\ge x\right]\le e^{-2x^2/T},\qquad
 \Prob\!\left[\max_{i,t\le T}(-N_{i,t})\ge x\right]
 \le\min\{1,ne^{-2x^2/T}\}.
\]
For completeness, the nonnegative supermartingale
$\exp(\lambda M_t-\lambda^2t/8)$ gives the first bound after
optimizing $\lambda$; apply the same argument to $-N_{i,t}$
and union over $i$ for the second.

Deterministically,
\[
 R_t\le\sum_{s\le t}(\mu_{I_s}-\mu_*)+M_t+\max_i(-N_{i,t}).
\]
The pseudo-regret increments are nonnegative, so their prefix is at
most $W$. Integrating the first tail gives a bound $\sqrt T$.
For the second, split at $a=\sqrt{T\log(2n)/2}$ and use
\[
 \int_a^\infty e^{-2x^2/T}\,dx
 \le\frac{T}{4a}e^{-2a^2/T}.
\]
Its expected maximum is at most
$a+T/(8a)\le2\sqrt{T\log(2n)}$. With~\eqref{eq:pseudoreg},
\begin{equation}\label{eq:maxprefix}
 \E\max_{0\le t\le T}R_t\le U+3\sqrt{T\log(2n)}=A.
\end{equation}
The random variable on the left is nonnegative because it includes
$t=0$. This controls the whole prefix path, not merely terminal
regret or pseudo-regret.

\subsection{A fallback preserving all-input regret}

The actual learner follows the pure learner until the end of the first
round $\tau$ with $R_\tau\ge\kappa=8A/\rho$. If $\tau<T$,
it starts fresh ordinary Hedge on the remaining rounds, with uniform
weights and learning rate
$\eta_H=\sqrt{8\log n/(T-\tau)}$. Its sampling tape is an independent
part of the original shared tape. No fallback is needed if $\tau=T$
or no crossing occurs.

The pure process is defined counterfactually for all $T$ rounds.
Before the first crossing it equals the actual process. Under any iid
law, Markov and~\eqref{eq:maxprefix} give
$\Prob(\tau\le T)\le\rho/8$. If paired pure paths agree and neither
run crosses, the actual paths agree. Thus actual path mismatch is at
most $\rho/8+2(\rho/8)=3\rho/8<\rho$.
This union bound makes no assumption about independence of the crossing
events.

On every fixed deterministic input,
$R_t-R_{t-1}\in[-1,1]$, since the selected cost and the increment of
the minimum cumulative expert cost each lie in $[0,1]$.
The prefix regret at the first crossing is therefore less than
$\kappa+1$. If there is no crossing, terminal regret is below $\kappa$.
Conditional on $\tau$ and the pure history, the suffix is a fixed
cost sequence and the fallback tape retains its independent law.

For a fixed suffix of length $h$, the elementary Hedge potential
calculation gives expected regret at most
$\log n/\eta+\eta h/8$: normalize exponential weights each round,
apply the bounded-variable log-moment inequality
$\log\E e^{-\eta X}\le-\eta\E X+\eta^2/8$, and compare the final
potential with the contribution of a best suffix expert.
With the stated $\eta_H$, this is
$\sqrt{h\log n/2}\le\sqrt{T\log n/2}$.
The sum of the best prefix and best suffix costs is at most the
best \emph{single full-horizon} expert's cost. Adding the two regret
bounds, then averaging over the complete original tape, proves
\eqref{eq:upper}. The minimum with $T$ uses the pathwise bound
$\Reg_T\le T$, not nonnegativity of realized regret.

The implementation stores cumulative costs and the current expert,
computes each race in $O(n)$ operations, monitors regret in $O(n)$
per round, and uses ordinary $O(n)$ Hedge updates after switching.
The counterfactual pure continuation is for analysis only and need
not be computed after switching. This proves the complexity
claim in Theorem~\ref{thm:upper}.

\section{What the separation does and does not imply}\label{sec:discussion}

The lower-bound family imposes sequential causal commitments to
independent phase biases. It does not restrict the tape to independent
per-phase coins or impose a per-phase replication budget.
The proof pools all randomness and charges disjoint first mismatches.
The cube of experts is essential to comparing every phase with
\emph{one} fixed full-horizon comparator.

Replicable coin estimation~\cite{HIKLY} permits adaptive choices of
which coordinate to sample. Its Section 6.4 and Lemma 6.14 also permit
sequential subroutines, aggregation of their transcript, and a rare
sample-cutoff fallback to a nonreplicable procedure. Thus sequential
composition, shared randomness and such a fallback are not new here.
Those results allow sample allocation and stopping rules that are not
freely available in the prescribed chronological phase stream. Their
coordinate-sampling lower bound (Theorem 4.46 and Appendix A.4) also
does not by itself impose the cumulative causal commitments used in
Theorem~\ref{thm:lower}.
The realizable PAC lower bound in~\cite{LMPS} concerns a final
classifier and an iid sample-complexity criterion. Matching exponents
of $\log n$ in those statements alone do not imply
Theorem~\ref{thm:lower}; a reduction preserving causality, the full
tape and cumulative regret would be needed. This is a scoped
comparison of these theorems, not an exhaustive priority claim.

At fixed $n$ our lower bound remains of square-root horizon order.
It does not resolve the general optimal product-replicable rate.
The linear-versus-sublinear regime uses
$n=\exp(\Theta(T^{1/3}))$; the second regime has $n<T$ but gives
a multiplicative cost gap rather than linear regret.
Exact entire-path replication is stronger than agreement on a
final recommendation or low average action disagreement.
No trained-model degradation, numerical advantage, empirical benchmark,
external endorsement, or hiring consequence is established by the proofs.

\begin{thebibliography}{9}
\bibitem{ABB}
S. Ahmadi, S. Bhandari and A. Blum.
\emph{Replicable Online Learning}.
Advances in Neural Information Processing Systems, 2025.
\url{https://proceedings.neurips.cc/paper_files/paper/2025/file/da8abfd153886b745278b40b5df0be93-Paper-Conference.pdf}.

\bibitem{HIKLY}
M. Hopkins, R. Impagliazzo, D. Kane, S. Liu and C. Ye.
\emph{Replicability in High Dimensional Statistics}.
arXiv:2406.02628v1, 2024.
\url{https://arxiv.org/abs/2406.02628v1}.

\bibitem{LMPS}
K. G. Larsen, M. E. Mathiasen, C. Pabbaraju and C. Svendsen.
\emph{The Sample Complexity of Replicable Realizable PAC Learning}.
arXiv:2602.19552v1, 2026.
\url{https://arxiv.org/abs/2602.19552v1}.
\end{thebibliography}
\end{document}
